We have conducted the first systematic empirical test of the Modifiable Areal
Unit Problem in algorithmic redistricting.  Across a 130$\times$ unit-count
range---from county equivalents ($\sim$100,000 pop each) to census block
groups ($\sim$1,200 pop each)---the bisect recursive bisection algorithm
produces Polsby-Popper compactness scores that vary by less than 2\% for 48
of 50 states.\textsuperscript{\dag}  D-seat predictions are identical across
all three resolutions for all five focus states tested.\textsuperscript{\dag}

\paragraph{Structural explanation.}
MAUP-robustness in bisect arises from three reinforcing mechanisms: the
algorithm optimises a graph-edge-cut objective with no direct dependence on
node area; Polsby-Popper is insensitive to small boundary fluctuations; and
the population-balance constraint ($\pm 0.5\%$) eliminates most of the freedom
that would otherwise allow resolution-dependent variation.

\paragraph{Practical upshot.}
For the two states that exceed the 2\% threshold (Idaho, Wyoming, both $k=2$),
the cause is a small-$k$ effect rather than MAUP, and the absolute magnitude
of variation (RSI $\approx 2.3$--$2.6\%$) remains small in absolute PP terms
($\Delta PP \approx 0.005$).

\paragraph{Contribution.}
This paper validates a previously untested assumption underlying the bisect
portfolio: that using census tracts as the atomic unit does not materially
affect conclusions.  Researchers can cite this paper when defending tract-level
findings to audiences that question resolution choices.  Practitioners using
bisect at block-group resolution can expect behaviour consistent with the
published tract-level characterisation.

\paragraph{Future work.}
Block-level ($\sim$11 million nodes nationally) redistricting remains
computationally challenging but not infeasible on modern hardware.  A
50-state block-level run would extend the tested unit-count range from
$130\times$ to $\sim$3{,}500$\times$ and is a natural next step.  Testing
MAUP sensitivity for VRA-weighted and county-weighted modes is also
warranted, as those modes alter edge weights in ways that may interact with
resolution differently than the geographic mode tested here.
